\subsection{Метод максимального правдоподобия}

Метод состоит в том, что в качестве «наиболее правдоподобного» значения параметра берут значение \( \theta \), которое максимизирует вероятность получить при \( n \) опытах данную выборку \( \vec{X}=(X_1, \ldots ,X_n) \). Это значение параметра \( \theta \) зависит от выборки и является искомой оценкой.

Распределение Бернулли \( \textrm{B}_\theta \) является дискретным, поэтому введём \textit{функцию правдоподобия} для нашей выборки при фиксированном \( \theta \):
\[
f_\theta(\vec{X}) = \prod_{i=1}^n P(X_i) = \prod_{i=1}^n \theta^{X_i}(1−\theta)^{1−X_i}
\]

В дискретном случае при фиксированных \( \vec{x}=(x_1,\ldots,x_n)^T \) значение \( f_\theta(\vec{x}) \) равно вероятности, с которой выборка \( \vec{X}=(X_1,\ldots,X_n)^T \) в данной серии экспериментов принимает значения \( x_1,\ldots,x_n \):
\[
    \begin{aligned}
        f_\theta(\vec{x}) &= P_\theta(X_1=x_1, \ldots, X_n=x_n)= \\
        &= P_\theta(X_1=x_1) \cdot \ldots \cdot P_\theta(X_n=x_n)
    \end{aligned}
\]

Отметим следующие наблюдения:
\begin{itemize}
    \item Генеральная совокупность \( X_i=\{0,1\},\ i=\overline{1,n} \).
    \item Экспоненциальная функция \( g_a(x)=a^x \) обладает свойством:
    \[
        g_a(x+y)=g_a(x)\cdot g_a(y)
    \]
\end{itemize}

На основе этих пунктов упростим функцию правдоподобия до вида:
\[
   f_\theta(\vec{X}) = \theta^S (1−\theta)^{n−S},
\]
где \( S \) —-- сумма элементов выборки.

На основе полученной функции получим искомую оценку.

\textbf{Оценкой максимального правдоподобия} \textit{(ОМП)} \( \hat{\theta} \) для неизвестного параметра \( \theta \) называют такое значение \( \theta \), при котором достигается максимум функции \( f_\theta(\vec{X}) \).

Для упрощения вычислений можно ввести \textit{логарифмическую функцию правдоподобия}:
\[
    L_\theta(\vec{X}) \equiv \ln f_\theta(\vec{X})= S\ln \theta + (n-S)\ln(1-\theta)
\]

Поскольку функция \( h(x)=\ln x \) монотонна, точки максимума функций \( f_\theta(\vec{X}) \) и \( L_\theta(\vec{X}) \) совпадают. Поэтому для оценки максимально правдоподобного \( \theta \) можно использовать любую из них.

Воспользуемся необходимым условием экстремума функции одной переменной \textit{(относительно \( \theta \))}:
\[
    \frac{dL}{d\theta} = 0\implies\frac{S}{\theta} - \frac{n-S}{1-\theta} = 0 \implies \hat{\theta} = \frac{S}{n} = \dfrac{1}{n}\sum_{i=1}^nX_i=\overline{X}
\]

Убедимся, что стационарная точка $\hat{\theta}$ --- точка глобального максимума. Вычислим вторую производную логарифмической функции:
\[
    \dfrac{d^2L}{d\theta^2} = -\frac{S}{\theta^2} - \frac{n - S}{(1 - \theta)^2}
\]

Для распределения Бернулли $S \geq 0$ и $n-S \geq 0$. Поэтому, в силу положительности знаменателей,
\[
    \forall\theta\in(0,1)\ \dfrac{d^2L}{d\theta^2}<0.
\]

Следовательно, функция $L_\theta$ строго выпукла вверх на $(0,1)$, а точка $\hat{\theta} = \bar{X}$ является единственной точкой глобального максимума.

Таким образом, оценка максимального правдоподобия для параметра $\theta$ распределения Бернулли совпадает со средним выборочным.

Вычислим ОМП при помощи скрипта:

\begin{lstlisting}[caption=Расчёт ОМП]
bernoulli = [
  0, 0, 0, 1, 1, 0, 0, 1, 1, 1,
  1, 1, 0, 1, 1, 0, 0, 1, 0, 1,
  0, 1, 1, 1, 1, 1, 1, 0, 0, 1,
  0, 1, 0, 0, 0, 0, 0, 1, 1, 1,
  1, 0, 1, 1, 1, 0, 0, 1, 0, 1,
]

n = len(bernoulli)
S = sum(bernoulli)
theta_hat = S / n

print(f"theta = {S} / {n} = {theta_hat}")
\end{lstlisting}

Таким образом,
\[
\hat{\theta} = \dfrac{28}{50} = 0.56
\]

Некоторые свойства полученной оценки:

\begin{enumerate}
    \item \textit{Несмещённость.}
    \[
        \begin{gathered}
            E[\hat{\theta}] = \dfrac{1}{n}\sum_{i=1}^nE[X_i]=\dfrac{1}{n}\cdot n\theta = \theta\\ E[\hat{\theta}]=\theta\ \square
        \end{gathered}
    \]
    \item \textit{Состоятельность.}
    \[
    \begin{gathered}
        \hat{\theta} = \overline{X} \xrightarrow[n \to \infty]{P} E[X_1] = \theta\\
        \hat{\theta}\xrightarrow[n\to\infty]{P}\theta\ \square
    \end{gathered}
    \]
    % \item \textit{Эффективность.} Оценка достигает нижней границы Рао--Крамера. Для распределения Бернулли информация Фишера равна $I(\theta) = \frac{1}{\theta(1-\theta)}$, а дисперсия оценки $D[\hat{\theta}] = \frac{\theta(1-\theta)}{n} = \frac{1}{n I(\theta)}$, что подтверждает её оптимальность в классе несмещённых оценок.
\end{enumerate}